\documentclass[11pt]{article}
\usepackage[T1]{fontenc}
\usepackage{textcomp}
\usepackage[margin=0.75in]{geometry}
\usepackage{amsmath,amssymb,amsthm}
\usepackage{array,booktabs}
\usepackage{hyperref}
\setlength{\parindent}{0pt}
\newtheorem{theorem}{Theorem}
\theoremstyle{definition}
\newtheorem{definition}{Definition}
\title{Flow Proof Artifact: Nat-sq-succ}
\author{Generated by \texttt{flow doc proof}}
\date{Flow Proof Book}
\begin{document}
\maketitle
Squaring a successor expands by self-multiplication.\\[0.5em]
\textbf{Source.} peano — https://en.wikipedia.org/wiki/Square\_(algebra)\\[0.5em]
\bigskip
\noindent{\large \textbf{Derived fact 1.} \textit{sq(succ(n)) = succ(n) * succ(n)} \hfill the natural numbers \cdot square \cdot successor squares by self multiplication}\par
\smallskip\noindent\textit{Coordinate: the natural numbers \cdot square \cdot successor squares by self multiplication}\par
\smallskip\noindent\textit{Source: peano}\par
\smallskip\noindent\textit{Built on: squaring is self-multiplication, for square on the natural numbers, successor on the right via commutativity, for addition on the natural numbers}\par
\medskip
\noindent\textbf{Goal.} sq(succ(n)) = succ(n) * succ(n)\par
\noindent$\displaystyle \forall n \in \mathbb{N}\quad sq(s) = s^{2}$\par
\medskip
\noindent
\begin{tabular}{@{} >{\raggedright\arraybackslash}p{0.06\textwidth} >{\raggedright\arraybackslash}p{0.47\textwidth} >{\raggedleft\arraybackslash}p{0.41\textwidth} @{}}
\textbf{\#} & \textbf{Proof} & \textbf{Mathematics} \\
\midrule
\textbf{1.} & We prove that successor squares by self multiplication for square on the natural numbers. & \\[0.45em]
\textbf{2.} & Let s = succ(n). & \\[0.45em]
\textbf{3.} & We invoke the definitional clause governing square on the natural numbers: squaring is self-multiplication, for square on the natural numbers (instantiated for s). & \\[0.45em]
\textbf{4.} & We invoke the derived fact governing addition on the natural numbers: successor on the right via commutativity, for addition on the natural numbers (instantiated for n, 0). & \\[0.45em]
\textbf{5.} & From step 2, step 3, and step 4, this implies sq(s) equals s times s. Hence proven. & $\displaystyle sq(s) = s^{2}$ \\[0.45em]
\bottomrule
\end{tabular}
\label{thm:Nat-square-successor_squares_by_self_multiplication}
\par\medskip\hrule
\end{document}
